The results support a clear but disciplined interpretation. CMWSL does not overturn the well-known fact that daily stock-index volatility is strongly persistence-driven. HAR remains the right benchmark, and in the main Rogers--Satchell specification it remains the strongest average model. That finding should not be minimized. It is precisely because the benchmark is strong that the conditional gains documented in this paper matter. A credible contribution in this literature is not universal dominance over HAR, but a precise account of when richer representation begins to add value.

The first substantive implication concerns forecast horizon. The results consistently show that the relative value of multiscale representation rises as the task moves from one-day forecasting toward five- and ten-day horizons. This pattern has a natural economic interpretation. At the shortest horizon, the next day's volatility is heavily determined by immediate persistence and local market continuation. At longer horizons, however, medium-scale adjustment processes, information diffusion, and broader market-risk conditions have more time to matter. The wavelet layer is useful precisely because it organizes those medium-scale dynamics without discarding the strong persistence structure already captured by HAR.

The second implication concerns information richness. CMWSL performs more strongly when the predictor set is expanded with public credit-spread, policy-uncertainty, and volatility-term-structure variables, and when the target is changed from Rogers--Satchell to Parkinson volatility. These are not incidental robustness wins. Together they suggest that the multiscale layer becomes more valuable when the forecasting environment contains more heterogeneous information or when the target places more weight on informative OHLC range variation. In other words, the wavelet design is not acting as a generic complexity booster; it is acting as an information organizer.

The third implication concerns the nature of cross-market dependence. The spillover results indicate that peer-market information is not uniformly useful. Its contribution is concentrated in broad-market indices and in stressed states, and it appears most clearly through the d2 and d3 wavelet bands. This point matters because it refines the usual spillover narrative. Existing work often treats cross-market dependence as a broad time-domain phenomenon. The present evidence suggests a more specific view: a meaningful part of risk transmission across indices occurs through medium- and long-scale channels associated with weekly to monthly reallocation, gradual risk repricing, and synchronized stress adjustment. That interpretation is especially compelling for DJIA and S\&P~500, and much less so for the technology-concentrated Nasdaq-100.

The divergence between DM and CW in the spillover comparison fits this interpretation. DM is conservative in nested settings because the larger model pays an estimation-noise penalty that is not explicitly corrected. CW is designed for exactly this situation. The fact that CW detects significant incremental spillover content where positive QLIKE gains are also observed implies that the peer-index wavelet block contains real predictive information, but that the signal is moderate rather than overwhelming. This combination strengthens rather than weakens the paper's claims: the evidence is robust enough to survive a correctly specified nested-model test, while remaining within the range that is realistic for financial time-series data of this kind.

The warning results lead to a related practical conclusion. The best warning model in the paper is not the most elaborate one. A simpler logistic classifier built on the same causal feature pipeline delivers the most stable precision--recall performance and the clearest operating trade-off between earlier detection and alert burden. This is important for market-risk practice. Warning systems are often judged not only by discrimination, but also by whether they can be audited, calibrated, and implemented by human decision makers. The logistic model is therefore a feature rather than a limitation of the paper's design.

\subsection{Implications for Financial Risk Management Practice}
\label{subsec:risk_management_implications}

The results carry several direct implications for market risk management practice, beyond their contribution to the academic forecasting literature.

First, the state-dependent pattern of CMWSL gains has immediate relevance for risk model governance. The finding that multiscale and cross-index improvements concentrate in upper-volatility regimes and synchronized stress windows implies that practitioners can adopt CMWSL as a selective overlay: using the HAR baseline under normal market conditions and switching to the full multiscale pipeline when macro-financial indicators or rolling volatility diagnostics signal an elevated-stress environment. This regime-conditional deployment is consistent with existing internal model validation frameworks that require documented conditions under which supplementary models become operative.

Second, the risk early warning classifier delivers an operationally tractable alert mechanism. The precision--recall trade-off documented across settings allows a risk manager or chief risk officer to calibrate the alert threshold to match the institution's tolerance for missed warnings versus false positives. Unlike black-box neural network warning systems, the logistic classifier built on CMWSL features is directly interpretable: feature coefficients indicate which wavelet bands and macro-financial variables are currently driving the alert probability, which supports both model validation and regulatory communication. This property is especially relevant under risk management frameworks that require human-intelligible model outputs.

Third, the public-data design strengthens the case for deployment. All predictors---including the Rogers--Satchell and Parkinson volatility targets, wavelet features, and macro-financial enrichment variables---are derived from publicly available daily market and FRED data. This means that the pipeline can be reproduced, back-tested, and independently verified without access to proprietary intraday feeds. For institutions subject to model risk management guidelines (such as the Basel Committee's principles for sound model risk management, or the U.S.\ SR~11-7 guidance), a transparent, reproducible, public-data pipeline with documented walk-forward evaluation records is directly aligned with validation requirements.

Fourth, the spillover evidence has portfolio-level implications. The finding that S\&P~500 volatility benefits most from DJIA and Nasdaq wavelet spillover at the ten-day horizon suggests that cross-index risk transmission at medium frequencies is a material input to multi-asset portfolio risk estimation. Portfolio risk managers monitoring equity correlations may find that adding medium-scale wavelet energy from peer indices to their risk factor models improves variance forecast accuracy during stress windows---which is precisely when diversification benefits are most likely to break down and portfolio Value-at-Risk estimates are most likely to underestimate actual exposure.

\subsection{What the Paper Contributes to the Literature}
\label{subsec:discussion_spillover}

Taken together, the findings support three contributions to the volatility-forecasting and risk-monitoring literature.

First, the paper contributes a stricter evaluation standard. By taking HAR seriously, using public data only, and enforcing chronology through the entire feature pipeline, it reduces the scope for inflated claims. This matters because many apparent machine-learning improvements in financial forecasting disappear once the benchmark and evaluation protocol become more demanding.

Second, the paper contributes a more specific understanding of multiscale information. The wavelet block is not helpful everywhere, but it becomes valuable in exactly the settings where a risk manager would expect linear persistence to become less sufficient: longer horizons, richer information environments, and stressed states. That conditional structure is a more useful contribution than a marginal full-sample average improvement would have been.

Third, the paper provides direct evidence on cross-index spillover at specific wavelet frequencies. The S\&P~500 \(h{=}10\) cell, and the broader DJIA and S\&P~500 CW pattern, confirm that peer-index d2/d3 wavelet energy adds predictive content beyond the target index's own history. This is the study's strongest novel result: it isolates a specific transmission channel rather than merely asserting that peer markets matter.

\subsection{Limitations and Next Steps}
\label{subsec:limitations}

The study has several limitations. The analysis is restricted to three major U.S.\ equity indices, so the spillover findings should not be over-generalized to other markets or asset classes. The public-data design improves reproducibility, but it excludes proprietary intraday information that might sharpen very short-horizon forecasting or clarify cross-market lead--lag effects. Some extended macro variables are not treated with full vintage-level release timing in the main benchmark, so the expanded-data setting should be interpreted as a demanding robustness design rather than as the cleanest causal baseline. In addition, the robustness exercises do not yet use the full 2016--2025 evaluation length of the main specification.

These limitations point directly to future work. The most important extension is to test the same spillover design on international index panels and cross-asset systems, where asynchronous trading hours and stronger lead--lag channels could make frequency-specific spillover even more informative. A second extension is to lengthen the robustness windows so that the richer-information and Parkinson-target experiments match the full main evaluation period. A third is to complement the current cell-by-cell DM and CW analysis with model-confidence-set procedures. Finally, more specialized recurrent or Transformer architectures could be evaluated, but only if they are held to the same standards of chronology, calibration, and public-data reproducibility used in the present paper.

A fourth direction concerns the integration of unstructured information and human-in-the-loop model validation. The current CMWSL pipeline relies entirely on structured public time-series data. Extending it to incorporate signals extracted from financial news, central bank communications, or macroeconomic disclosures would require robust structured-extraction pipelines with appropriate domain co-design. Recent work on co-designing AI-assisted information extraction with domain experts \cite{li2026endoextractcodesigningstructuredtext} and on characterizing the complementary error profiles of human annotators and large language models in structured extraction tasks \cite{li2026whofailswhere} offers methodological templates from adjacent domains that could inform the design of human-AI collaborative pipelines for financial risk signal extraction. Similarly, uncertainty-aware feature selection strategies \cite{li2025ougsactiveviewselection} could provide principled protocols for active model recalibration when the distribution of incoming market observations diverges from the rolling training window---a practically important scenario for real-time risk monitoring under structural breaks.
